\section{Conclusion}

Inspired by memory schema theory that highlights \emph{prefrontal regions} and \emph{hippocampus regions}, we present \ourmodel, a two-stage optimization framework that decouples \textbf{how to organize memory} from \textbf{what to store}. Specifically, \ourmodel first induces a transferable, schema-consistent guideline for memory evolution, and then optimizes a guideline-aligned memory policy to decide what to retain, update, or forget across multi-session interactions. Extensive experiments on three personalization memory benchmarks show that \ourmodel consistently outperforms strong retrieval-based memory bank, and RL-based memory-agent baselines, while remaining robust under longer histories and noisier evidence. Overall, the results support that coupling an explicit evolution guideline with policy optimization yields a practical improvement in \textbf{efficiency}, \textbf{robustness}, and \textbf{transferability} for evolving user memory in conversational agents.

\section*{Acknowledgements}
This work was supported in part by the grants from National Science and Technology Major Project (No. 2023ZD0121104), National Natural Science Foundation of China (No. U22B2059), the Anhui Natural Science Foundation (No. 2508085ZD006), National Natural Science Foundation of China (No.62502404), Hong Kong Research Grants Council (Research Impact Fund No.R1015-23, Collaborative Research Fund No.C1043-24GF, General Research Fund No. 11218325), Institute of Digital Medicine of City University of Hong Kong (No.9229503), Huawei (Huawei Innovation Research Program), Tencent (Tencent Rhino-Bird Focused Research Program, Tencent University Cooperation Project), Didi (CCF-Didi Gaia Scholars Research Fund), Kuaishou (CCF-Kuaishou Large Model Explorer Fund No. 2025008, Kuaishou University Cooperation Project), and Bytedance.


\section*{Limitations}
Overall, our method is effective for improving long-horizon personalization memory by learning for more structured and consistent memory evolution. However, the second-stage optimization relies on an LLM-based scorer to provide guideline-aligned process rewards, which makes performance sensitive to scorer reliability. Moreover, our method requires careful tuning of the per-round token budget and the number of evolution rounds; when long histories are split into many rounds, small update errors can compound over time and lead to unintended forgetting or over-generalized memory entries. Finally, our current design treats memory evolution as a single-objective policy under a fixed guideline; extending it to explicitly balance multiple competing objectives (e.g., stability vs.\ plasticity, informativeness vs.\ brevity) remains non-trivial and may require additional control mechanisms.

\section*{Ethical considerations}
Our method is a general memory-evolution framework intended to support personalized agents, and it primarily improves \emph{how} existing memories are organized and optimized rather than expanding the scope of the LLM system's access. The primary ethical risk arises from misuse rather than from the method itself: if deployed without appropriate safeguards, persistent memory could be used to over-collect user information or to enable unwanted profiling. In responsible deployments, memory should follow data-minimization principles, avoid storing sensitive identifiers, and provide clear user controls for inspection, correction, and deletion; additionally, retention policies and access control should be enforced at the system level to ensure the system remains aligned with privacy expectations as application contexts evolve.

